\section{Sistemi che utilizzano LLM e RAG per rispondere in maniera automatica} \label{sec:sistemi-simili}

I chatbot per l'assistenza universitaria non nascono con i LLM: le prime generazioni di questi sistemi, risalenti a
ben prima della diffusione dei Transformer, si basavano su alberi di decisione scritti a mano o su classificatori
statistici, abbinati a un insieme limitato e predefinito di risposte. Erano sistemi capaci di gestire bene le
domande più frequenti, previste in fase di progettazione, ma fragili di fronte a qualunque formulazione che si
discostasse dal copione atteso. L'arrivo di LLM in grado di comprendere una domanda espressa liberamente, combinati
con il paradigma RAG per ancorare la risposta a documenti reali invece che a un insieme fisso di risposte
prestabilite, segna il passaggio a una nuova generazione di sistemi: capaci, in linea di principio, di rispondere a
qualunque domanda per cui esista un documento pertinente nella base indicizzata, senza dover prevedere a priori ogni
possibile formulazione.

L'applicazione di LLM e RAG all'assistenza universitaria non è, a questo punto, un'idea isolata. Peyton et al.
\cite{peyton2025review}, in una rassegna recente dedicata ai chatbot universitari per il supporto agli studenti,
confrontano otto implementazioni distinte e osservano come il settore non abbia ancora sfruttato appieno il
potenziale di questi sistemi rispetto ad altri ambiti applicativi, in parte per la difficoltà di reperire e curare
dataset di dominio chiuso, specifici di un singolo ateneo, su cui valutarli in modo affidabile. È un'osservazione
che motiva direttamente la scelta di questo lavoro di costruire un dataset diagnostico dedicato
(Capitolo~\ref{cap:validazione}), invece di affidarsi a benchmark generici.

Tre sistemi recenti sono particolarmente vicini, per dominio e impostazione, a quello sviluppato in questo lavoro.
Unipa-GPT \cite{siragusa2024unipagpt}, sviluppato per l'Università di Palermo, è un assistente conversazionale in
italiano rivolto a potenziali studenti che devono scegliere un corso di laurea triennale o magistrale. Gli autori
confrontano esplicitamente un approccio RAG, che recupera il contesto da un corpus raccolto dal sito web
dell'ateneo, con un approccio di fine-tuning dello stesso modello di base: un punto di riferimento diretto, su un
dominio analogo a quello di questo progetto, per la stessa alternativa architetturale discussa nella
Sezione~\ref{sec:rag}. Marcel \cite{trienes2025marcel}, sviluppato per l'Università di Marburg, si concentra invece
sulle domande più frequenti dei futuri studenti relative all'ammissione, e affianca alla ricerca sui documenti un
retriever dedicato alle FAQ: in questo modo il personale amministrativo può orientare esplicitamente le risposte
alle domande ricorrenti, invece di affidarsi interamente all'indicizzazione automatica dei documenti originali.
Vlahek \cite{vlahek2024chatbot}, infine, descrive un prototipo di chatbot RAG per il supporto agli studenti
universitari realizzato con uno strumento \emph{low-code} (Flowise AI), invece che con una pipeline sviluppata su
misura come nei due casi precedenti. Il lavoro conferma che la combinazione di LLM e RAG per questo tipo di compito
è ormai accessibile anche senza un'implementazione interamente personalizzata, benché a scapito del controllo
puntuale su ciascuno stadio della pipeline (chunking, ricerca ibrida, reranking) discusso nella
Sezione~\ref{sec:rag}, centrale nelle scelte di questo lavoro.

Il quadro non è però limitato all'ambito universitario: sistemi LLM più RAG vengono ormai distribuiti anche in
contesti industriali, e la letteratura che ne documenta la realizzazione è complementare a quella dei sistemi
accademici, perché si concentra sui problemi di ingegnerizzazione che emergono una volta messo in uso un sistema
reale, più che sulla sola qualità delle risposte misurata su un benchmark. Barnett et al.
\cite{barnett2024sevenfailures} riportano l'esperienza di tre implementazioni RAG distinte, in ambito di ricerca,
educativo e biomedico, e ne ricavano sette punti di fallimento ricorrenti. Tra questi il \emph{missing content}, in
cui il documento necessario a rispondere non è presente nella base indicizzata, e il \emph{missed top-ranked
document}, in cui il documento pertinente esiste nell'indice ma non emerge tra i candidati restituiti dal recupero.
Gli autori osservano inoltre che la robustezza di un sistema RAG si costruisce iterativamente con l'uso reale e non
può essere interamente garantita in fase di progettazione. Bourdin et al. \cite{bourdin2025easirag} documentano un
caso concreto di questo tipo in un contesto manifatturiero: un laboratorio di test ambientali che ha realizzato un
assistente RAG per rispondere alle domande degli operatori a partire dalle procedure operative interne.

Letti insieme, i sistemi universitari descritti condividono con questo lavoro l'impostazione di fondo --- LLM più
RAG, applicato alle domande degli studenti di un singolo ateneo --- ma nessuno descrive come gestire una domanda la
cui risposta dipende dal corso di laurea dello studente, quando questa informazione non è esplicita nel testo della
domanda. Unipa-GPT e Vlahek operano su un unico corpus (rispettivamente l'intero ateneo di Palermo e un singolo
corso), dove il problema non si pone nella stessa forma; Marcel, pur limitato alle sole domande di ammissione, non
discute casi di ambiguità tra profili di studenti diversi.

È in questo spazio che si colloca il presente lavoro. Il sistema sviluppato indicizza insieme le pagine di più corsi
di laurea, le pagine generali di ateneo e i materiali didattici: la stessa domanda --- una scadenza, un requisito,
un referente --- può quindi avere risposte diverse a seconda del corso. Per questo il sistema riconosce quando una
domanda è ambigua rispetto al corso, chiede allo studente di indicarlo, e usa poi il corso noto per privilegiare i
documenti pertinenti, senza escludere quelli degli altri corsi o di ateneo. La soluzione adottata è descritta nel
Capitolo~\ref{cap:architettura} e valutata nel Capitolo~\ref{cap:validazione}.
